\hypertarget{appendix-f1a6d43e80d1}{}
\section*{Screen-space ray generation and mouse picking}
\section{Screen-space ray generation and mouse picking}\label{screen-space-ray-generation-and-mouse-picking}

\textbf{Curriculum depth:} Developed · \textbf{First seen:} Phase 11 · \textbf{Latest development:} Phase 11

\subsection{Definition and mechanism}\label{definition-and-mechanism}

A screen ray begins with a pixel location mapped through the viewport to NDC, then through inverse projection and camera pose into a world-space origin and direction. For perspective cameras, rays usually share the eye origin; orthographic rays have parallel directions but different origins across the image. Mouse picking must use the same high-DPI scaling, viewport offset, origin, and Y convention as rendering. Near and far unprojection can define a ray segment, but a point at infinity or an infinite far plane needs an explicit direction construction. Picking then intersects that ray with scene geometry or uses an object-ID buffer.

\subsection{Prerequisites}\label{prerequisites}

\begin{itemize}
\tightlist
\item
  \href{}{Camera view and clip-space pipeline (04-Visual-Computing-II, appendix-1c07d64f5e70)}
\item
  \hyperref[appendix-f744e41b95b4]{Unit directions and safe normalization}
\item
  \href{}{Visual angle and calibrated rays (04-Visual-Computing-II, appendix-7d356c396d1a)}
\end{itemize}

\subsection{Used by}\label{used-by}

\begin{itemize}
\tightlist
\item
  No downstream concept node is added in this stage.
\end{itemize}

\subsection{Full teaching explanation}\label{full-teaching-explanation}

\begin{itemize}
\tightlist
\item
  \href{Continuous\%20Notes/Computing\%20Foundations\%20-\%20From\%20Bits\%20to\%20C++\%20and\%20Systems\%20APIs\#Recover\%20geometry\%20from\%20depth\%20and\%20construct\%20screen-space\%20rays}{Recover geometry from depth and construct screen-space rays}
\end{itemize}

\subsection{Source coverage}\label{source-coverage}

Supplied lesson parts: 1082--1083.

\subsection{Future keywords}\label{future-keywords}

inverse projection, viewport mapping, picking ray
